\subsection{EQUIVALENCE OF MATRICES OVER A FIELD}

PROPOSITION 13. Let E, F be two finite-dimensional vector spaces over a field K. If \(u: \mathbf{E} \to \mathbf{F}\) is a linear mapping of rank \(r\) , there exist bases of E and F such that, with respect to these bases,

\[
M (u) = \left( \begin{array}{c c} I _ {r} & 0 \\ 0 & 0 \end{array} \right).\tag{55}
\]

Every matrix of type \((m, n)\) over \(\mathbf{K}\) and of rank \(r\) is equivalent to a matrix of the form (55).

The second assertion is trivially equivalent to the first. To show the latter, let \(\dim \mathbf{E} = n\) , \(\dim \mathbf{F} = m\) . The kernel \(\mathbf{N} = \bar{u}^{1}(0)\) is of dimension \(n - r\) ( \(\S 7\) , no. 4, formula (11)); let \(\mathbf{V}\) be a supplementary subspace of \(\mathbf{N}\) in \(\mathbf{E}\) and \((e_i)_{1 \leqslant i \leqslant n}\) a basis of \(\mathbf{E}\) such that \((e_i)_{1 \leqslant i \leqslant r}\) is a basis of \(\mathbf{V}\) and \((e_i)_{r + 1 \leqslant i \leqslant n}\) a basis of \(\mathbf{N}\) . Then the \(u(e_j)\) ( \(1 \leqslant j \leqslant r\) ) form a basis of \(u(\mathbf{E})\) ; hence there exists a basis \((f_j)_{1 \leqslant j \leqslant m}\) of \(\mathbf{F}\) such that \(f_j = u(e_j)\) for \(1 \leqslant j \leqslant r\) ( \(\S 7\) , no. 1, Theorem 2)

and clearly with respect to the bases \((e_{i})\) and \((f_{j})\) the matrix \(M(u)\) is given by (55).

COROLLARY. For two matrices over a field, of type \((m, n)\) , to be equivalent, it is necessary and sufficient that they have the same rank.

We shall now recover Proposition 13 by another more explicit method. For every ring A, every \(\lambda \in \mathbf{A}\) , every integer \(m > 1\) and every ordered pair of distinct integers \(i, j\) in \([1, m]\) , we write

\[
B _ {i j} (\lambda) = I _ {m} + \lambda E _ {i j}\tag{56}
\]

an invertible matrix of order \(m\) by no. 8.

Lemma 1. Let \(X = (\xi_{ij})\) be a matrix of type \((m, n)\) over a ring \(A\) . Suppose that \(m \geqslant 2\) and that there exists an element \(\xi_{i1}\) in the first column of \(X\) which is invertible in \(A\) . Then there exist two invertible square matrices \(P \in \mathbf{M}_m(A)\) , \(Q \in \mathbf{M}_n(A)\) and a matrix \(Y\) of type \((m - 1, n - 1)\) over \(A\) such that \(P\) (resp. \(Q\) ) is a product of matrices of the form \(B_{ij}(\lambda)\) of order \(m\) (resp. \(n\) ) and

\[
P X Q = \left( \begin{array}{c c c c} 1 & 0 & \ldots & 0 \\ 0 & & \\ \vdots & & Y \\ 0 & & \end{array} \right).\tag{57}
\]

The matrix \(B_{ij}(\lambda)X\) is obtained by adding to the row of \(X\) of index \(i\) the row of index \(j\) multiplied on the left by \(\lambda\) (no. 9, Example 2); if \(\xi_{i1}\) is invertible, then there exists \(\lambda \in A\) such that, for the matrix \(X' = B_{1i}(\lambda)X = (\xi_{kl}')\) , \(\xi_{11}' = 1\) ; multiplying \(X'\) on the left by suitably chosen matrices \(B_{k1}(\mu_k)\) of order \(m\) (for \(1 \leqslant k \leqslant m\) ), a matrix \(X'' = (\xi_{kl}')\) is obtained such that \(\xi_{j1}' = 1\) , \(\xi_{k1}'' = 0\) for \(k \neq 1\) . Then the matrix obtained is multiplied successively on the right by suitable matrices \(B_{1j}(\nu_j)\) of order \(n\) ( \(2 \leqslant j \leqslant n\) ) and a matrix is obtained of the form (57).

PROPOSITION 14. Let X be a matrix of type \((m, n)\) over a field K. If X is of rank r, there exist two invertible square matrices \(P \in \mathbf{M}_{m}(\mathbf{K})\) , \(Q \in M_{n}(\mathbf{K})\) such that P (resp. Q) is a product of matrices of order m (resp. n) of the form \(B_{ij}(\lambda)\) and

\[
P X Q = \left( \begin{array}{c c c c c c c} 1 & 0 & \dots & 0 & 0 & \dots & 0 \\ 0 & 1 & \dots & 0 & 0 & \dots & 0 \\ \cdot & \cdot & \cdot & \cdot & \cdot & \cdot & \cdot \\ 0 & 0 & \dots & \delta_ {r} & 0 & \dots & 0 \\ 0 & 0 & \dots & 0 & 0 & \dots & 0 \\ \cdot & \cdot & \cdot & \cdot & \cdot & \cdot & \cdot \\ 0 & 0 & \dots & 0 & 0 & \dots & 0 \end{array} \right)\tag{58}
\]

(a matrix \((\eta_{ij})\) all of whose terms are zero except the \(\eta_{ii}\) for \(1 \leqslant i \leqslant r\) , with \(\eta_{ii} = 1\) for \(1 \leqslant i \leqslant r - 1\) , \(\eta_{rr} = \delta_r \neq 0\) ). If \(r \neq m\) or \(r \neq n\) , it may also be assumed that \(\delta_r = 1\) .

The proposition is obvious if X = 0; suppose therefore \(X \neq 0\) . If m = n = 1 the proposition is obvious (with \(P = I_{m}\) , \(Q = I_{n}\) , \(\delta_{1} \neq 0\) arbitrary). If n = 1, \(m \geqslant 2\) , we can apply Lemma 1 (since \(X \neq 0\) ), which gives the desired form (58) with r = 1, \(\delta_{r} = 1\) . We argue by induction on n \textgreater{} 1; there exists an element \(\xi_{ij} \neq 0\) in X; if j = 1, Lemma 1 can be applied and reduces the problem to the case where X has the form (57). The induction hypothesis then applies to Y and there are therefore invertible matrices

\[
P ^ {\prime} \in \mathbf {M} _ {m - 1} (\mathbf {K}), \quad Q ^ {\prime} \in \mathbf {M} _ {n - 1} (\mathbf {K})
\]

which are products of matrices of the form \(B_{ij}(\lambda)\) of order m - 1 (resp. n - 1), such that \(P^{\prime}YQ^{\prime}\) is of the form (58). But, if \(B_{ij}(\lambda)\) belongs for example to \(\mathbf{M}_{m-1}(\mathbf{K})\) , then

\[
\left( \begin{array}{c c} 1 & 0 \\ 0 & B _ {i j} (\lambda) \end{array} \right) = B _ {i + 1, j + 1} (\lambda);
\]

formula (58) then follows from the formula for block products writing

\[
P = \left( \begin{array}{c c} 1 & 0 \\ 0 & P ^ {\prime} \end{array} \right) \quad \text { and } \quad Q = \left( \begin{array}{c c} 1 & 0 \\ 0 & Q ^ {\prime} \end{array} \right).
\]

If finally \(j \neq 1\) , it would be sufficient to consider the matrix \(XB_{j1}(1)\) to reduce it to the above case.

Proposition 14 recovers Proposition 13 immediately.

COROLLARY 1. If \(X\) is an invertible square matrix of order \(n\) over a field \(\mathbf{K}\) , there exist three invertible matrices \(P, Q, D\) of order \(n\) such that \(X = PDQ\) , \(P\) and \(Q\) being products of matrices of the form \(B_{ij}(\lambda)\) and \(D\) a diagonal matrix of the form

\[
D = \operatorname{diag} (1, 1, \dots , \delta),
\]

where \(\delta \neq 0\) (cf.~Exercise 13).

COROLLARY 2. For every field K, the group of invertible matrices \(\mathbf{GL}(n,\mathbf{K})\) is generated by the permutation matrices (no. 7, Example 2), the diagonal matrices \(\text{diag}(a,1,\ldots,1)\) ( \(a\neq0\) in K) and the matrices \(B_{12}(\lambda)\) ( \(\lambda\in K\) ).

It has been seen (no. 9) that the right (resp. left) product of a matrix by the matrix of a suitable transposition exchanges any two columns (resp. rows). Then the matrix \(\text{diag}(1, \ldots, 1, a)\) is equal to the product of \(\text{diag}(a, 1, \ldots, 1)\) and permutation matrices and every matrix \(B_{ij}(\lambda)\) is equal to the product of \(B_{12}(\lambda)\) and permutation matrices, whence the corollary.

Remarks. (1) In Chapter III, we shall see that, if \(m = n = r\) and \(K\) is commutative, then, for all choices of \(P\) and \(Q\) satisfying the conditions of Proposition 14, the element \(\delta_r\) is always the same and equal to the determinant of \(X\) (III, § 8, no. 6).

\begin{enumerate}
\def\labelenumi{(\arabic{enumi})}
\setcounter{enumi}{1}
\tightlist
\item
  The argument of Proposition 14, slightly modified, shows that there is a permutation matrix \(R\) such that (with the same conditions on \(P\) )
\end{enumerate}

\[
P X R = \left( \begin{array}{c c} I _ {r} & N \\ 0 & 0 \end{array} \right)
\]

if \(m = n = r\) does not hold, and

\[
P X R = \operatorname{diag} (1, \dots , 1, \delta)
\]

otherwise. Observe also that the method of proof gives an explicit determination of the matrices P, Q, R when X is given explicitly.

